\section{Modelo de Procesos Web (UWE), Secuencias y Máquinas de Estados}
\label{sec:modelo_procesos_uwe}
\label{sec:diagramas_secuencia}
\label{sec:diagramas_estados}
% =============================================================================

En la metodología UWE (Koch et al., 2008), el \textbf{Modelo de Procesos} refina el comportamiento interno de cada nodo estereotipado como \texttt{«process class»} en el Modelo de Navegación mediante un flujo de trabajo estructurado (Diagrama de Actividades de Proceso UWE) y lo vincula con la dinámica temporal de intercambio de mensajes (Diagramas de Secuencia UML) y el ciclo de vida de los artefactos (Diagramas de Máquina de Estados Finitos UML).

\subsection{Diagrama de Actividades del Proceso UWE: ``Emitir y Custodiar Voto''}
\label{subsec:actividades_proceso_uwe}

La Figura \ref{fig:uml_actividades} especifica el Diagrama de Actividades del proceso central \texttt{«process class» ProcesoEmisionVoto}, detallando la secuencia algorítmica completa desde la confirmación de la opción en cabina hasta la entrega del recibo firmado: (1) selección de candidato, (2) empaquetado de longitud fija a 32 bytes (1 byte de opción + 31 bytes CSPRNG), (3) encapsulamiento de clave simétrica mediante ML-KEM-768 contra $ek_{\text{AC}}$, (4) cifrado autenticado AES-256-GCM vinculando $\text{AAD} = \langle id_{\text{elec}}, ct_{\text{kem}} \rangle$, (5) firma digital ML-DSA-65 con contexto \texttt{ELEC\_BOLETA\_V1}, (6) transmisión REST (\texttt{POST /api/votar}), (7) verificación multinivel y desacoplamiento transaccional en el Servidor de Votación, y (8) emisión del recibo SHA3-256 firmado con $sk_{\text{SV}}$.

\begin{figure}[!htbp]
	\centering
	\resizebox{!}{0.68\textheight}{%
	\begin{tikzpicture}[
		scale=0.9,
		act/.style={rectangle, draw=black, fill=white, line width=0.8pt, rounded corners=1.5pt, font=\scriptsize, align=center, minimum height=0.72cm, text width=5.0cm},
		dec/.style={diamond, draw=black, fill=white, line width=0.8pt, aspect=2.1, font=\tiny, align=center, inner sep=1.5pt},
		errbox/.style={rectangle, draw=black, fill=white, double, double distance=1.2pt, line width=0.7pt, rounded corners=1.5pt, font=\scriptsize, align=center, text width=3.1cm},
		flow/.style={-{Latex[length=2mm, width=1.4mm]}, line width=0.8pt, draw=black},
		init/.style={circle, fill=black, inner sep=3.5pt},
		stop/.style={circle, draw=black, line width=0.8pt, fill=white, inner sep=3.5pt}
	]
		% Nodo Inicial
		\node[init] (start) at (0, 0) {};
		
		% Actividades Cliente (Proceso UWE en Cabina)
		\node[act, below=0.45cm of start] (a1) {1. Confirmar opción en cabina y empaquetar voto fijo\\$m_{\text{fijo}} \leftarrow \text{Pad}_{32\text{B}}(\text{opción} \;\|\; \text{CSPRNG}_{31\text{B}})$};
		\node[act, below=0.45cm of a1] (a2) {2. Encapsular clave efímera contra clave pública de AC\\$(K, ct_{\text{kem}}) \leftarrow \text{ML-KEM-768.Encaps}(ek_{\text{AC}})$};
		\node[act, below=0.45cm of a2] (a3) {3. Cifrar con AES-256-GCM ($\text{AAD}=\langle id_{\text{elec}}, ct_{\text{kem}} \rangle$) y firmar\\$c_{\text{voto}} \leftarrow \text{AEAD.Enc}(K, m_{\text{fijo}}, \text{AAD})$; \; $\sigma_V \leftarrow \text{ML-DSA-65.Sign}(sk_V)$};
		\node[act, below=0.45cm of a3] (a4) {4. Transmitir \texttt{PaqueteVoto} por API REST (\texttt{POST /api/votar})};
		
		% Decisiones Servidor
		\node[dec, below=0.55cm of a4] (d1) {¿Credencial AR válida\\(\texttt{ELEC\_CRED\_AR\_V1})?};
		\node[dec, below=0.75cm of d1] (d2) {¿Firma del votante $\sigma_V$\\válida (\texttt{ELEC\_BOLETA\_V1})?};
		\node[dec, below=0.75cm of d2] (d3) {¿Registro en padrón\\indica \texttt{ha\_votado = true}?};
		
		% Rechazos
		\node[errbox, right=2.1cm of d1] (err1) {Abortar proceso UWE:\\Credencial no autorizada};
		\node[errbox, right=2.1cm of d2] (err2) {Abortar proceso UWE:\\Firma digital corrupta};
		\node[errbox, right=2.1cm of d3] (err3) {Disparar CU9: Rechazar\\intento de doble voto};
		
		% Actividades Éxito
		\node[act, below=0.75cm of d3] (a5) {5. Transacción atómica SQL: \texttt{ha\_votado = TRUE} y desacoplar\\identidad depositando tupla anónima en \texttt{tbl\_anfora\_boletas}};
		\node[act, below=0.45cm of a5] (a6) {6. Calcular $h_b = \text{SHA3-256}(\text{Boleta})$ y firmar recibo con $sk_{\text{SV}}$\\(\texttt{ctx = "ELEC\_RECIBO\_SV\_V1"})};
		\node[act, below=0.45cm of a6] (a7) {7. Activar restricción OCL \texttt{NoRetornoTrasVotar} y renderizar\\vista \texttt{PaginaReciboInmutable}};
		
		% Nodo Final
		\node[stop, below=0.45cm of a7] (final) {};
		\node[circle, fill=black, inner sep=2pt] at (final.center) {};
		
		% Conectores
		\draw[flow] (start) -- (a1);
		\draw[flow] (a1) -- (a2);
		\draw[flow] (a2) -- (a3);
		\draw[flow] (a3) -- (a4);
		\draw[flow] (a4) -- (d1);
		
		\draw[flow] (d1) -- node[left, font=\tiny\itshape] {Sí} (d2);
		\draw[flow] (d1) -- node[above, font=\tiny\itshape] {No} (err1);
		
		\draw[flow] (d2) -- node[left, font=\tiny\itshape] {Sí} (d3);
		\draw[flow] (d2) -- node[above, font=\tiny\itshape] {No} (err2);
		
		\draw[flow] (d3) -- node[left, font=\tiny\itshape] {No} (a5);
		\draw[flow] (d3) -- node[above, font=\tiny\itshape] {Sí} (err3);
		
		\draw[flow] (a5) -- (a6);
		\draw[flow] (a6) -- (a7);
		\draw[flow] (a7) -- (final);
		
		% Retornos de error al final
		\node[stop, right=6.0cm of d2] (err_final) {};
		\node[circle, fill=black, inner sep=2pt] at (err_final.center) {};
		\draw[flow] (err1) -| (err_final);
		\draw[flow] (err2) -- (err_final);
		\draw[flow] (err3) -| (err_final);
	\end{tikzpicture}
	}%
	\caption{Diagrama de Actividades del Proceso UWE (\texttt{ProcesoEmisionVoto}) para la emisión, validación y custodia disociada del voto.}
	\label{fig:uml_actividades}
	\par\vspace{3pt}{\footnotesize\textit{Fuente: Elaboración propia.}}
\end{figure}

\noindent\textbf{Declaración de Tensión Arquitectónica (Criptografía en Cliente vs.\ Servidor — \texttt{PEND-15}):}
Es metodológicamente fundamental explicitar la relación entre este Modelo de Procesos UWE y la implementación demostrativa actual:
\begin{itemize}
	\item \textbf{Especificación Objetivo en el Modelo de Procesos UWE:} Los pasos 1, 2 y 3 de la Figura \ref{fig:uml_actividades} (empaquetado de 32 bytes, encapsulamiento ML-KEM-768 y firma ML-DSA-65 con $sk_V$) se especifican para ejecutarse íntegramente en el entorno local del cliente web (terminal de cabina), de modo que el voto en texto claro y la clave privada $sk_V$ jamás transiten hacia el servidor.
	\item \textbf{Estado Actual en el Prototipo Web Flask (\texttt{PEND-15}):} Dado que el motor criptográfico nativo (\texttt{libpqc\_core.dylib}) se enlaza mediante Python FFI (\texttt{ctypes}), el módulo demostrativo \texttt{src/web/app.py} recibe la selección desde el navegador y ejecuta la invocación de encapsulamiento y firma con \texttt{\_sk\_votante} en el proceso de backend antes de entregar el paquete al módulo de urna. Esta divergencia se documenta con total transparencia en la Matriz de Conformidad (Tabla \ref{tab:conformidad_prototipo_diseno}, ítem \texttt{PEND-15}), estableciendo que el modelo UWE constituye la especificación formal invariante hacia la cual converge la evolución del software (mediante compilación de \texttt{PQClean} a WebAssembly en el navegador).
\end{itemize}

\subsection{Acoplamiento con los Diagramas de Secuencia UML del Protocolo}

El Modelo de Procesos UWE se proyecta directamente sobre la línea temporal de mensajes entre los actores del protocolo en dos secuencias operativas: la fase de emisión, validación multinivel y custodia del voto (Figura \ref{fig:uml_secuencia_votacion}), y la fase de clausura comicial, permutación de lote (\textit{shuffle}), desencapsulamiento en la estación \textit{air-gapped} de la AC y publicación del acta auditable (Figura \ref{fig:uml_secuencia_conteo}).

\begin{figure}[!htbp]
	\centering
	\resizebox{\textwidth}{!}{%
	\begin{tikzpicture}[
		node distance=1.1cm and 3.2cm,
		lifeline/.style={rectangle, draw=black, fill=white, line width=0.8pt, rounded corners=1.5pt, minimum width=2.8cm, minimum height=0.7cm, font=\scriptsize, align=center},
		msg/.style={-{Latex[length=2mm, width=1.4mm]}, line width=0.8pt, draw=black},
		retmsg/.style={-{Latex[length=2mm, width=1.4mm]}, line width=0.7pt, dashed, draw=black},
		note/.style={rectangle, draw=black, fill=white, line width=0.6pt, rounded corners=1pt, font=\tiny, align=left, inner sep=3.5pt}
	]
		% Cabeceras
		\node[lifeline] (V)  at (0, 0) {Votante ($V$)};
		\node[lifeline] (AR) at (5.0, 0) {Autoridad Reg. ($AR$)};
		\node[lifeline] (SV) at (10.2, 0) {Servidor Votación ($SV$)};
		
		% Ejes verticales
		\draw[draw=black, dashed, line width=0.5pt] (V) -- ++(0,-9.6);
		\draw[draw=black, dashed, line width=0.5pt] (AR) -- ++(0,-9.6);
		\draw[draw=black, dashed, line width=0.5pt] (SV) -- ++(0,-9.6);
		
		% Bloques de activación
		\draw[fill=white, draw=black, line width=0.6pt] ([xshift=-1.5mm, yshift=-0.8cm]V.center) rectangle ++(3mm, -1.8cm);
		\draw[fill=white, draw=black, line width=0.6pt] ([xshift=-1.5mm, yshift=-1.2cm]AR.center) rectangle ++(3mm, -1.2cm);
		\draw[fill=white, draw=black, line width=0.6pt] ([xshift=-1.5mm, yshift=-3.4cm]V.center) rectangle ++(3mm, -3.2cm);
		\draw[fill=white, draw=black, line width=0.6pt] ([xshift=-1.5mm, yshift=-4.8cm]SV.center) rectangle ++(3mm, -3.8cm);
		
		% Mensajes Registro
		\draw[msg] ([yshift=-1.2cm]V.center) -- ([yshift=-1.2cm]AR.center)
		node[midway, above, font=\scriptsize] {1: \texttt{SolicitudRegistro(id\_v, pk\_v)}};
		
		\draw[retmsg] ([yshift=-2.2cm]AR.center) -- ([yshift=-2.2cm]V.center)
		node[midway, above, font=\scriptsize] {2: \texttt{CredencialVotante(id\_v, pk\_v, ts, firma\_ar)}};
		
		% Cifrado local
		\node[note, right=0.2cm of V, yshift=-3.8cm] {
			\textbf{Proceso UWE en Cabina (\texttt{ProcesoEmisionVoto}):}\\
			$\bullet$ Cargar $ek_{\text{AC}}$ pública de la Autoridad de Conteo\\
			$\bullet$ $(K, ct_{kem}) \leftarrow \text{ML-KEM.Encaps}(ek_{\text{AC}})$\\
			$\bullet$ $c_{voto} \leftarrow \text{AES-GCM.Encrypt}(K, m_{\text{fijo}}, nonce, \text{AAD}=\langle id_{\text{elec}}, ct_{kem} \rangle)$\\
			$\bullet$ $\text{Boleta} = \langle id_{\text{elec}}, ct_{kem}, c_{voto}, nonce \rangle$\\
			$\bullet$ $\sigma_V \leftarrow \text{ML-DSA.Sign}(sk_V, \text{canonical}(\text{Boleta}))$\\
			$\bullet$ $\text{PaqueteVoto} = \langle \text{Boleta}, \sigma_V, \text{Credencial} \rangle$
		};
		
		% Envío de voto
		\draw[msg] ([yshift=-5.2cm]V.center) -- ([yshift=-5.2cm]SV.center)
		node[midway, above, font=\scriptsize] {3: \texttt{POST /api/votar} $(\text{PaqueteVoto})$};
		
		% Validación en Servidor
		\node[note, left=0.2cm of SV, yshift=-6.2cm] {
			\textbf{Verificación Multinivel y Disociación en SV:}\\
			$\bullet$ $\text{ML-DSA.Verify}(\sigma_{AR}, \text{Credencial}, pk_{AR}) \stackrel{?}{=} \text{OK}$\\
			$\bullet$ $\text{ML-DSA.Verify}(\sigma_V, \text{Boleta}, pk_V) \stackrel{?}{=} \text{OK}$\\
			$\bullet$ Comprobación y bloqueo atómico: \texttt{ha\_votado = FALSE}\\
			$\bullet$ Transacción atómica en BD: \texttt{Padron.marcarYaVoto()} e\\
			\quad \texttt{Urna.insert(Boleta)} (sin credencial ni firma $\sigma_V$)\\
			$\bullet$ $\text{Recibo} = \langle h(\text{Boleta}), ts_{\text{ventana}} \rangle$; \; $\sigma_{SV} \leftarrow \text{ML-DSA.Sign}(sk_{SV}, \text{Recibo})$
		};
		
		% Retorno de Recibo
		\draw[retmsg] ([yshift=-8.2cm]SV.center) -- ([yshift=-8.2cm]V.center)
		node[midway, above, font=\scriptsize] {4: \texttt{ReciboVotacion(h\_boleta, ts\_ventana, firma\_sv)}};
	\end{tikzpicture}
	}%
	\caption{Diagrama de Secuencia UML: Emisión desde el proceso UWE, verificación multinivel y retorno de recibo firmado.}
	\label{fig:uml_secuencia_votacion}
	\par\vspace{3pt}{\footnotesize\textit{Fuente: Elaboración propia.}}
\end{figure}

\begin{figure}[!htbp]
	\centering
	\resizebox{\textwidth}{!}{%
	\begin{tikzpicture}[
		node distance=1.1cm and 3.5cm,
		lifeline/.style={rectangle, draw=black, fill=white, line width=0.8pt, rounded corners=1.5pt, minimum width=2.8cm, minimum height=0.7cm, font=\scriptsize, align=center},
		msg/.style={-{Latex[length=2mm, width=1.4mm]}, line width=0.8pt, draw=black},
		retmsg/.style={-{Latex[length=2mm, width=1.4mm]}, line width=0.7pt, dashed, draw=black},
		note/.style={rectangle, draw=black, fill=white, line width=0.6pt, rounded corners=1pt, font=\tiny, align=left, inner sep=3.5pt}
	]
		% Cabeceras
		\node[lifeline] (SV) at (0, 0) {Servidor Votación ($SV$)};
		\node[lifeline] (AC) at (5.5, 0) {Autoridad Conteo ($AC$)};
		\node[lifeline] (PUB) at (11.0, 0) {Portal Auditoría UWE};
		
		% Ejes verticales
		\draw[draw=black, dashed, line width=0.5pt] (SV) -- ++(0,-8.2);
		\draw[draw=black, dashed, line width=0.5pt] (AC) -- ++(0,-8.2);
		\draw[draw=black, dashed, line width=0.5pt] (PUB) -- ++(0,-8.2);
		
		% Activaciones
		\draw[fill=white, draw=black, line width=0.6pt] ([xshift=-1.5mm, yshift=-0.8cm]SV.center) rectangle ++(3mm, -1.8cm);
		\draw[fill=white, draw=black, line width=0.6pt] ([xshift=-1.5mm, yshift=-1.8cm]AC.center) rectangle ++(3mm, -5.2cm);
		\draw[fill=white, draw=black, line width=0.6pt] ([xshift=-1.5mm, yshift=-6.4cm]PUB.center) rectangle ++(3mm, -1.2cm);
		
		% Mensajes
		\draw[msg] ([yshift=-1.4cm]SV.center) -- ([yshift=-1.4cm]AC.center)
		node[midway, above, font=\scriptsize] {1: Clausura de mesa $\to$ Entrega de Lote Permutado (Shuffle)};
		\node[note, below=0.1cm of SV, xshift=1.2cm, yshift=-1.6cm] {
			\textbf{Permutación en SV:}\\
			Aplica ordenación pseudoaleatoria al lote anónimo\\
			$\{\langle id_{\text{elec}}, ct_{kem}, c_{voto}, nonce \rangle_j\}$ destruyendo orden temporal.
		};
		
		% Nota Desencapsulamiento
		\node[note, right=0.2cm of AC, yshift=-3.5cm] {
			\textbf{Procesamiento por la AC (Estación Air-Gapped):}\\
			$\bullet$ Carga de clave privada de desencapsulamiento $dk_{\text{AC}}$ (resguardada en HSM)\\
			$\bullet$ \textbf{Para cada boleta anónima} $j \in \{1, \dots, N\}$:\\
			\quad 1) $K_j \leftarrow \text{ML-KEM.Decaps}(dk_{\text{AC}}, ct_{kem, j})$\\
			\quad 2) $m_j \leftarrow \text{AES-GCM.Decrypt}(K_j, c_{voto, j}, nonce_j, \text{AAD}=\langle id_{\text{elec}}, ct_{kem, j} \rangle)$\\
			\quad 3) Desempaquetar voto de longitud fija (32 B) y validar opción\\
			\quad 4) Acumular cómputo: $\text{Total}[m_j] \leftarrow \text{Total}[m_j] + 1$\\
			\quad 5) Registrar digest inmutable: $h_j = \text{SHA3-256}(\text{Boleta}_j)$
		};
		
		% Publicación
		\draw[msg] ([yshift=-6.6cm]AC.center) -- ([yshift=-6.6cm]PUB.center)
		node[midway, above, font=\scriptsize] {2: Publicar Acta Oficial de Cómputo + Lista de Hashes $\{h_j\}$};
		
		\node[note, left=0.2cm of PUB, yshift=-7.4cm] {
			\textbf{Auditoría Ciudadana en Portal UWE (\texttt{«query»}):}\\
			Los votantes verifican la presencia del hash de su recibo en la lista pública oficial.
		};
	\end{tikzpicture}
	}%
	\caption{Diagrama de Secuencia UML: Escrutinio por la Autoridad de Conteo y publicación en el Portal de Auditoría UWE.}
	\label{fig:uml_secuencia_conteo}
	\par\vspace{3pt}{\footnotesize\textit{Fuente: Elaboración propia.}}
\end{figure}

\subsection{Acoplamiento con las Máquinas de Estados Finitos UML}

Las transiciones entre los nodos de navegación y procesos UWE guardan correspondencia biunívoca con las dos máquinas de estados finitos del sistema: el ciclo de vida de la boleta de sufragio (Figura \ref{fig:uml_estados_boleta}) y los estados institucionales del proceso electoral (Figura \ref{fig:uml_estados_eleccion}).

\begin{figure}[!htbp]
	\centering
	\begin{tikzpicture}[
		scale=0.88,
		state/.style={rectangle, draw=black, fill=white, line width=0.8pt, rounded corners=1.5pt, minimum height=0.75cm, font=\scriptsize, align=center, inner sep=4pt},
		errstate/.style={rectangle, draw=black, fill=white, double, double distance=1.2pt, line width=0.7pt, rounded corners=1.5pt, minimum height=0.75cm, font=\scriptsize, align=center, inner sep=4pt},
		trans/.style={-{Latex[length=2mm, width=1.4mm]}, line width=0.7pt, draw=black, font=\tiny\itshape},
		init/.style={circle, fill=black, inner sep=3.5pt},
		final/.style={circle, draw=black, line width=0.8pt, fill=white, inner sep=3.5pt}
	]
		\node[init] (s0) at (0, 0) {};
		\node[state] (s1) at (2.4, 0) {1. En Edición\\(Cabina UWE)};
		\node[state] (s2) at (6.6, 0) {2. Cifrada Híbrida\\(KEM-DEM)};
		\node[state] (s3) at (10.8, 0) {3. Firmada Digitalmente\\(ML-DSA-65)};
		
		\node[state] (s4) at (10.8, -2.5) {4. En Transmisión\\(Canal Red Dolev--Yao)};
		\node[state] (s5) at (6.6, -2.5) {5. Custodiada en Urna\\(Ánfora Disociada)};
		\node[state] (s6) at (2.4, -2.5) {6. Desencapsulada\\(Apertura con $dk_{\text{AC}}$)};
		
		\node[final] (sf) at (0, -2.5) {};
		\node[circle, fill=black, inner sep=2pt] at (sf.center) {};
		
		% Estados de Error
		\node[errstate] (err_rej) at (8.7, -4.6) {Rechazada\\(Firma Inválida / No Empadronado)};
		\node[errstate] (err_dup) at (4.0, -4.6) {Descartada\\(Intento de Voto Duplicado)};
		
		% Transiciones de Éxito
		\draw[trans] (s0) -- (s1) node[midway, above] {Emitir};
		\draw[trans] (s1) -- (s2) node[midway, above] {$\text{Encaps}(ek_{\text{AC}})$};
		\draw[trans] (s2) -- (s3) node[midway, above] {$\text{Sign}(sk_V)$};
		\draw[trans] (s3) -- (s4) node[midway, right] {Enviar};
		\draw[trans] (s4) -- (s5) node[midway, above] {Validación OK};
		\draw[trans] (s5) -- (s6) node[midway, above] {Escrutinio};
		\draw[trans] (s6) -- (sf) node[midway, above] {Totalizar};
		
		% Transiciones de Falla
		\draw[trans, dashed] (s4) |- (err_rej) node[near start, right] {Firma corrupta};
		\draw[trans, dashed] (s4) |- (err_dup) node[near start, left] {Ya votó};
	\end{tikzpicture}
	\caption{Diagrama de Estados UML del ciclo de vida de la boleta de sufragio acoplado al proceso UWE.}
	\label{fig:uml_estados_boleta}
	\par\vspace{3pt}{\footnotesize\textit{Fuente: Elaboración propia.}}
\end{figure}

\begin{figure}[!htbp]
	\centering
	\begin{tikzpicture}[
		scale=0.9,
		st/.style={rectangle, draw=black, fill=white, line width=0.8pt, rounded corners=1.5pt, font=\scriptsize, align=center, inner sep=5pt, minimum height=0.72cm},
		arr/.style={-{Latex[length=2mm, width=1.4mm]}, line width=0.7pt, draw=black, font=\tiny\itshape},
		i/.style={circle, fill=black, inner sep=3.5pt},
		f/.style={circle, draw=black, line width=0.8pt, fill=white, inner sep=3.5pt}
	]
		\node[i] (start) at (0, 0) {};
		\node[st] (st1) at (2.4, 0) {1. Configuración\\y Ceremonia Claves};
		\node[st] (st2) at (6.8, 0) {2. Empadronamiento\\y Registro AR};
		\node[st] (st3) at (11.2, 0) {3. Jornada Electoral\\(Votación Abierta)};
		
		\node[st] (st4) at (11.2, -2.3) {4. Clausura Comicial\\(Recepción Cerrada)};
		\node[st] (st5) at (6.8, -2.3) {5. Desencapsulamiento\\y Escrutinio AC};
		\node[st] (st6) at (2.4, -2.3) {6. Auditoría UWE y\\Proclamación Oficial};
		\node[f] (end) at (0, -2.3) {};
		\node[circle, fill=black, inner sep=2pt] at (end.center) {};
		
		\draw[arr] (start) -- (st1);
		\draw[arr] (st1) -- node[midway, above] {Claves listas} (st2);
		\draw[arr] (st2) -- node[midway, above] {Apertura comicios} (st3);
		\draw[arr] (st3) -- node[midway, right] {Cierre de mesas} (st4);
		\draw[arr] (st4) -- node[midway, above] {Lote permutado} (st5);
		\draw[arr] (st5) -- node[midway, above] {Publicar acta} (st6);
		\draw[arr] (st6) -- (end);
	\end{tikzpicture}
	\caption{Diagrama de Estados UML del proceso comicial institucional.}
	\label{fig:uml_estados_eleccion}
	\par\vspace{3pt}{\footnotesize\textit{Fuente: Elaboración propia.}}
\end{figure}
